%% =========================================================
%% Chapter: Replica entropy metrics
%% The common label function of the replica metrics and its marked ratios
%% (area-law paper, Section 6.2, Lemma 6.2, replicas:metric).
%% =========================================================

\chapter{Replica Entropy Metrics}\label{ch:replica_metrics}

The symmetric-replica argument of the two-dimensional area
law~\cite{OpenAI2026AreaLaw} compares the label observable of a subsystem with
a positive operator built from one label function shared by every subsystem.
Fix the dimension $d\ge1$ of the whole one-copy system and $0<t<1/4$. Pad each
label $\lambda\vdash k$ to length $d$, put $l_i=\lambda_i+d-i$ for
$1\le i\le d$, and set
\[
    c=d+(1+t)\binom d2,\qquad
    w_k(\lambda)=\frac{\Gamma(c)}{\Gamma(c+tk)}
        \prod_{i=1}^d\frac{\Gamma(1+tl_i)}{\Gamma(1+t(d-i))}.
\]
The metric of a subsystem $Q$ is $W_{Q,k}=\sum_\lambda w_k(\lambda)\pi^\lambda_{Q,k}$.
This chapter proves the scalar properties of $w_k$ on which the operator
statements of the paper's Lemma~6.2 rest: positivity, the exact marked ratio,
the entropy asymptotics, and the size and limit of the marked ratio. The
integral representation of $W_{Q,k}$ and the operator forms of these estimates
are not yet formalized.

\section{Ratios of Gamma values}

\begin{lemma}[Wendel's inequalities]
    \label{lem:wendel_gamma_ratio}
    \lean{Real.Gamma_add_le_rpow_mul_Gamma, Real.mul_Gamma_le_rpow_mul_Gamma_add,
        Real.Gamma_add_div_Gamma_le_rpow, Real.div_rpow_le_Gamma_add_div_Gamma,
        Real.rpow_sub_le_Gamma_add_div_Gamma}
    \leanok
    For $x>0$ and $0\le s\le1$,
    \[
        x(x+s)^{s-1}\le\frac{\Gamma(x+s)}{\Gamma(x)}\le x^s .
    \]
    If moreover $x\ge1$, then $x^s-s\le\Gamma(x+s)/\Gamma(x)$.
\end{lemma}

\begin{proof}\leanok
    Write $x+s=(1-s)x+s(x+1)$ and $x+1=s(x+s)+(1-s)(x+s+1)$, apply
    log-convexity of $\Gamma$ to each, and use $\Gamma(y+1)=y\Gamma(y)$. For
    the last claim, $(x+s)^{s-1}\ge x^s/(x+s)$, so the lower bound is at least
    $x^s-sx^s/(x+s)\ge x^s-s$ when $x\ge1$.
\end{proof}

\begin{lemma}[Stirling estimate for the Gamma function]
    \label{lem:log_gamma_stirling}
    \lean{Real.abs_log_Gamma_one_add_sub_le}
    \leanok
    \uses{lem:wendel_gamma_ratio, lem:stirling_log_factorial}
    For $u\ge0$,
    $|\log\Gamma(1+u)-(u\log u-u)|\le3\log(u+2)+3$.
\end{lemma}

\begin{proof}\leanok
    Let $n=\lfloor u\rfloor$ and $s=u-n$. Wendel's inequalities at $x=n+1$
    give $|\log\Gamma(1+u)-\log n!|\le\log(u+2)$. The Stirling bounds control
    $\log n!-(n\log n-n)$, and $u\log u-u$ differs from $n\log n-n$ by at most
    $\log(u+2)+2$.
\end{proof}

\section{The common label function}

\begin{lemma}[Positivity, normalization and marked ratio]
    \label{lem:replica_weight_ratio}
    \lean{Partition.replicaConst, Partition.replicaWeight, Partition.replicaRatio,
        Partition.replicaWeight_pos, Partition.replicaWeight_zero,
        Partition.replicaWeight_update_succ}
    \leanok
    The label function is positive, and $w_0=1$. If $\lambda\vdash k$ is
    obtained from $\nu\vdash k-1$ by adding a box in row $i$, then
    $w_k(\lambda)=r_{\lambda,i}\,w_{k-1}(\nu)$ with
    \[
        r_{\lambda,i}=\frac{\Gamma(c+t(k-1))}{\Gamma(c+tk)}
            \frac{\Gamma(1+tl_i)}{\Gamma(1+t(l_i-1))},
    \]
    where $l_i$ is the shifted part of $\lambda$.
\end{lemma}

\begin{proof}\leanok
    Every argument of $\Gamma$ is positive. Adding the box raises the number of
    copies by one and the $i$-th shifted part by one, leaving the other shifted
    parts unchanged.
\end{proof}

\begin{lemma}[Entropy asymptotics of the label function]
    \label{lem:replica_weight_entropy}
    \lean{Partition.exists_abs_log_replicaWeight_add_le}
    \leanok
    \uses{lem:log_gamma_stirling, lem:replica_weight_ratio}
    There is a constant $C_{d,t}$ such that every $\lambda\vdash k$ of length at
    most $d$ satisfies
    $|\log w_k(\lambda)+tkH(\lambda/k)|\le C_{d,t}\log(k+2)$.
\end{lemma}

\begin{proof}\leanok
    Apply the Stirling estimate to each of the $d+2$ Gamma factors. The terms
    containing $\log t$ cancel up to a constant because
    $\sum_il_i=k+\binom d2$. The shifts $0\le l_i-\lambda_i\le d-1$ change
    $\sum_il_i\log l_i$ by $O_d(\log(k+2))$, and
    $kH(\lambda/k)=k\log k-\sum_i\lambda_i\log\lambda_i$.
\end{proof}

Together with the entropy asymptotics of the dimension formula, this gives
$w_k(\lambda)=d_\lambda^{-t}$ up to factors polynomial in $k$, which is the
label-wise form of the comparison of $W_{Q,k}$ with $e^{-tF_{Q,k}}$.

\section{The marked ratio}

\begin{lemma}[Size of the marked ratio]
    \label{lem:replica_ratio_bounds}
    \lean{Partition.replicaRatio_eq_div, Partition.replicaRatio_wendel,
        Partition.exists_replicaRatio_bounds}
    \leanok
    \uses{lem:wendel_gamma_ratio, lem:replica_weight_ratio}
    Let $0<t<1$. There is $C\ge1$ such that, for $k\ge1$ and $1\le l\le k+d$,
    the ratio $r(k,l)$ of the previous lemma satisfies
    $C^{-1}(l/k)^t\le r(k,l)\le C$.
\end{lemma}

\begin{proof}\leanok
    The ratio is a quotient of two Wendel ratios, with $x=1+t(l-1)$ and
    $y=c+t(k-1)$, both at least one. Wendel's inequalities give
    $(1-t)(x/y)^t\le r(k,l)\le(x/y)^t/(1-t)$. On the stated range,
    $tl\le x\le(1+d)y$ and $y\le(c+t)k$.
\end{proof}

\begin{lemma}[Limit of the marked ratio]
    \label{lem:replica_ratio_limit}
    \lean{Real.abs_rpow_sub_rpow_le, Partition.exists_abs_replicaRatio_sub_le}
    \leanok
    \uses{lem:replica_ratio_bounds}
    Let $0<t<1$. There is a constant $C$ such that, for $k\ge1$ and
    $1\le l\le k+d$,
    \[
        \bigl|r(k,l)-\max\bigl((l-d)/k,0\bigr)^t\bigr|\le Ck^{-t}.
    \]
\end{lemma}

\begin{proof}\leanok
    Wendel's inequalities give $|r(k,l)-(x/y)^t|\le Ck^{-t}$, since $y\ge tk$
    and $(x/y)^t\le1+d$. Moreover $|x/y-l/k|=O(1/k)$ and
    $|l/k-\max((l-d)/k,0)|\le d/k$, and $|a^t-b^t|\le|a-b|^t$ for $a,b\ge0$.
\end{proof}

On a removable branch the star operator has eigenvalue $(\lambda_i-i)/k=(l_i-d)/k$,
so this lemma is the scalar content of the convergence of the marked ratio of
the metrics to the positive part of the star operator raised to the power $t$.
